\documentclass[10pt,a4paper]{amsart}
\usepackage{amsmath,amssymb,amsthm,graphicx}
% Retained preamble for the Stacks Project chapter divisors.tex (Divisors).
% The chapter file pinned-original-source.tex carries only \input{preamble} on its first line, so no
% preamble text lives inside the chapter itself.  The Stacks Project uses one single shared preamble for
% every chapter; the snapshot stages it verbatim in sources/stacks-chow/preamble.tex.  The block below is
% copied mechanically, line for line, from that file's lines 14-18 (the Xy-pic setup) and lines 154-217
% (the native theorem declarations together with the mathematical macros).  Nothing was added, removed or
% redefined.  The wrapper class and the amsmath/amssymb/amsthm/graphicx base packages are supplied by the
% assembler, so the class/\externaldocument part of the shared preamble is not repeated here.
% For commutative diagrams we use Xy-pic
\usepackage[all]{xy}

% We use 2cell for 2-commutative diagrams.
\xyoption{2cell}
% Theorem environments.
%
\theoremstyle{plain}
\newtheorem{theorem}[subsection]{Theorem}
\newtheorem{proposition}[subsection]{Proposition}
\newtheorem{lemma}[subsection]{Lemma}

\theoremstyle{definition}
\newtheorem{definition}[subsection]{Definition}
\newtheorem{example}[subsection]{Example}
\newtheorem{exercise}[subsection]{Exercise}
\newtheorem{situation}[subsection]{Situation}

\theoremstyle{remark}
\newtheorem{remark}[subsection]{Remark}
\newtheorem{remarks}[subsection]{Remarks}

\numberwithin{equation}{subsection}

% Macros
%
\def\lim{\mathop{\mathrm{lim}}\nolimits}
\def\colim{\mathop{\mathrm{colim}}\nolimits}
\def\Spec{\mathop{\mathrm{Spec}}}
\def\Hom{\mathop{\mathrm{Hom}}\nolimits}
\def\Ext{\mathop{\mathrm{Ext}}\nolimits}
\def\SheafHom{\mathop{\mathcal{H}\!\mathit{om}}\nolimits}
\def\SheafExt{\mathop{\mathcal{E}\!\mathit{xt}}\nolimits}
\def\Sch{\mathit{Sch}}
\def\Mor{\mathop{\mathrm{Mor}}\nolimits}
\def\Ob{\mathop{\mathrm{Ob}}\nolimits}
\def\Sh{\mathop{\mathit{Sh}}\nolimits}
\def\NL{\mathop{N\!L}\nolimits}
\def\CH{\mathop{\mathrm{CH}}\nolimits}
\def\proetale{{pro\text{-}\acute{e}tale}}
\def\etale{{\acute{e}tale}}
\def\QCoh{\mathit{QCoh}}
\def\Ker{\mathop{\mathrm{Ker}}}
\def\Im{\mathop{\mathrm{Im}}}
\def\Coker{\mathop{\mathrm{Coker}}}
\def\Coim{\mathop{\mathrm{Coim}}}

% Boxtimes
%
\DeclareMathSymbol{\boxtimes}{\mathbin}{AMSa}{"02}

%
% Macros for moduli stacks/spaces
%
\def\QCohstack{\mathcal{QC}\!\mathit{oh}}
\def\Cohstack{\mathcal{C}\!\mathit{oh}}
\def\Spacesstack{\mathcal{S}\!\mathit{paces}}
\def\Quotfunctor{\mathrm{Quot}}
\def\Hilbfunctor{\mathrm{Hilb}}
\def\Curvesstack{\mathcal{C}\!\mathit{urves}}
\def\Polarizedstack{\mathcal{P}\!\mathit{olarized}}
\def\Complexesstack{\mathcal{C}\!\mathit{omplexes}}
% \Pic is the operator that assigns to X its picard group, usage \Pic(X)
% \Picardstack_{X/B} denotes the Picard stack of X over B
% \Picardfunctor_{X/B} denotes the Picard functor of X over B
\def\Pic{\mathop{\mathrm{Pic}}\nolimits}
\def\Picardstack{\mathcal{P}\!\mathit{ic}}
\def\Picardfunctor{\mathrm{Pic}}
\def\Deformationcategory{\mathcal{D}\!\mathit{ef}}

% --- appended 2026-10-12 (Chief): the project own definitions of the reference/slogan/history
% display environments, copied mechanically from the shared Stacks preamble lines 8-12
% (\usepackage{verbatim} plus the three \newenvironment lines, which the chapter preamble keeps
% verbatim from the project preamble). Nothing added, removed or redefined: these environments
% carry no mathematical content and the project itself typesets them as comments. Without these
% three lines a counted statement whose own text uses them fails with Environment slogan undefined.
\usepackage{verbatim}
\newenvironment{reference}{\comment}{\endcomment}
\newenvironment{slogan}{\comment}{\endcomment}
\newenvironment{history}{\comment}{\endcomment}
\usepackage{cleveref}
\crefname{originalproposition}{original proposition}{original propositions}
\Crefname{originalproposition}{Original proposition}{Original propositions}
\makeatletter
\newlabel{algebra-proposition-dimension}{{00KQ}{source}{}{sourcepointer-1}{}}
\newlabel{algebra-proposition-dimension@cref}{{[originalproposition][0][]00KQ}{[1][0][]source}{}{}{}}
\makeatother
\crefname{originallemma}{original lemma}{original lemmas}
\Crefname{originallemma}{Original lemma}{Original lemmas}
\makeatletter
\newlabel{coherent-lemma-affine-if-quasi-affine}{{0EBE}{source}{}{sourcepointer-2}{}}
\newlabel{coherent-lemma-affine-if-quasi-affine@cref}{{[originallemma][0][]0EBE}{[1][0][]source}{}{}{}}
\makeatother
\crefname{originallemma}{original lemma}{original lemmas}
\Crefname{originallemma}{Original lemma}{Original lemmas}
\makeatletter
\newlabel{coherent-lemma-flat-base-change-cohomology}{{02KH}{source}{}{sourcepointer-3}{}}
\newlabel{coherent-lemma-flat-base-change-cohomology@cref}{{[originallemma][0][]02KH}{[1][0][]source}{}{}{}}
\makeatother
\crefname{originallemma}{original lemma}{original lemmas}
\Crefname{originallemma}{Original lemma}{Original lemmas}
\makeatletter
\newlabel{coherent-lemma-quasi-coherence-higher-direct-images}{{01XJ}{source}{}{sourcepointer-4}{}}
\newlabel{coherent-lemma-quasi-coherence-higher-direct-images@cref}{{[originallemma][0][]01XJ}{[1][0][]source}{}{}{}}
\makeatother
\crefname{originallemma}{original lemma}{original lemmas}
\Crefname{originallemma}{Original lemma}{Original lemmas}
\makeatletter
\newlabel{coherent-lemma-quasi-coherence-higher-direct-images-application}{{01XK}{source}{}{sourcepointer-5}{}}
\newlabel{coherent-lemma-quasi-coherence-higher-direct-images-application@cref}{{[originallemma][0][]01XK}{[1][0][]source}{}{}{}}
\makeatother
\crefname{originallemma}{original lemma}{original lemmas}
\Crefname{originallemma}{Original lemma}{Original lemmas}
\makeatletter
\newlabel{lemma-structure-sheaf-Xs}{{0EBN}{source}{}{sourcepointer-6}{}}
\newlabel{lemma-structure-sheaf-Xs@cref}{{[originallemma][0][]0EBN}{[1][0][]source}{}{}{}}
\makeatother
\crefname{originallemma}{original lemma}{original lemmas}
\Crefname{originallemma}{Original lemma}{Original lemmas}
\makeatletter
\newlabel{morphisms-lemma-characterize-affine}{{01S8}{source}{}{sourcepointer-7}{}}
\newlabel{morphisms-lemma-characterize-affine@cref}{{[originallemma][0][]01S8}{[1][0][]source}{}{}{}}
\makeatother
\crefname{originallemma}{original lemma}{original lemmas}
\Crefname{originallemma}{Original lemma}{Original lemmas}
\makeatletter
\newlabel{properties-lemma-characterize-affine}{{01QF}{source}{}{sourcepointer-8}{}}
\newlabel{properties-lemma-characterize-affine@cref}{{[originallemma][0][]01QF}{[1][0][]source}{}{}{}}
\makeatother
\title{When the inclusion of the locus where a section is invertible is affine}
\author{The Stacks Project Authors}
\date{stacks-project:divisors.tex:local-snapshot-20261012; GNU FDL 1.2 or later}
\begin{document}
\maketitle
\setcounter{subsection}{1}
\noindent\textit{Editorial note.} All delivered mathematical text is the author's own, copied from the pinned Stacks Project chapter file divisors.tex; no statement or proof step was rewritten, reordered or removed.\par\medskip
\noindent\textit{Editorial note.} Carrier window 3: statement and proof both lie inside the window 6000-8000 granted to this lane by the owner lane lane-s13 on 2026-10-14 (CONTROL 8.83 follow-up); lane-s15/carve.txt records the same edges.\par\medskip
\noindent\textit{Editorial note.} The prefilter flagged this candidate for 'the result is clear' in the cases x in U and dim 0; those are the trivial branches of the induction and the rest of the proof is argued in full, so the card is taken under the 2026-10-14 judge-rule.\par\medskip
\noindent\textit{Editorial note.} The proof contains one xymatrix square; it is preserved verbatim with the chapter's own Xy-pic setup.\par\medskip
\noindent\textit{Editorial note.} Preamble: sources/stacks-divisors/retained-preamble.tex (88 lines, verbatim present).\par\medskip
\noindent\textit{Allowed standard prerequisites (editorial).}\begin{itemize}
\item The structure sheaf of U as a colimit of the reflexive powers of F. Lemma lemma-structure-sheaf-Xs (7456-7476), used at proof lines 7569-7579 and 7620-7624. Turns elements of Gamma(U, O\_U) into sections of some F\textasciicircum{}\{[n]\}.
\item Affineness is local on the base. Morphisms, Lemma morphisms-lemma-characterize-affine, used at proof lines 7590-7592. Reduces the converse to X affine.
\item A quasi-affine scheme with vanishing higher cohomology is affine. Cohomology of Schemes, Lemma coherent-lemma-affine-if-quasi-affine, used at proof lines 7593-7595. Is the criterion applied at the end of the reduction.
\item Cohomology of the inclusion commutes with flat base change. Cohomology of Schemes, Lemmas coherent-lemma-quasi-coherence-higher-direct-images-application, coherent-lemma-quasi-coherence-higher-direct-images and coherent-lemma-flat-base-change-cohomology, used at proof lines 7595-7614. Reduces the vanishing to a local ring.
\item Characterisation of affine schemes by vanishing on a cover. Properties, Lemma properties-lemma-characterize-affine, used at proof lines 7628-7631. Concludes affineness of U from the affine pieces D(f\_i) cap U.
\end{itemize}\medskip
\ifdefined\hypertarget\hypertarget{sourcepointer-1}{}\fi
\noindent\textit{Source reference.} Original proposition 00KQ: focused-seeds/AGP/sources/stacks-snapshot-20261011/tags.txt, source lines 760--760. Navigation pointer into the Stacks Project; the standard fact it supplies is applied inside the counted proof span and the argument resting on it stays there.\par\medskip
\ifdefined\hypertarget\hypertarget{sourcepointer-2}{}\fi
\noindent\textit{Source reference.} Original lemma 0EBE: focused-seeds/AGP/sources/stacks-snapshot-20261011/tags.txt, source lines 17919--17919. Navigation pointer into the Stacks Project; the standard fact it supplies is applied inside the counted proof span and the argument resting on it stays there.\par\medskip
\ifdefined\hypertarget\hypertarget{sourcepointer-3}{}\fi
\noindent\textit{Source reference.} Original lemma 02KH: focused-seeds/AGP/sources/stacks-snapshot-20261011/tags.txt, source lines 3372--3372. Navigation pointer into the Stacks Project; the standard fact it supplies is applied inside the counted proof span and the argument resting on it stays there.\par\medskip
\ifdefined\hypertarget\hypertarget{sourcepointer-4}{}\fi
\noindent\textit{Source reference.} Original lemma 01XJ: focused-seeds/AGP/sources/stacks-snapshot-20261011/tags.txt, source lines 2520--2520. Navigation pointer into the Stacks Project; the standard fact it supplies is applied inside the counted proof span and the argument resting on it stays there.\par\medskip
\ifdefined\hypertarget\hypertarget{sourcepointer-5}{}\fi
\noindent\textit{Source reference.} Original lemma 01XK: focused-seeds/AGP/sources/stacks-snapshot-20261011/tags.txt, source lines 2521--2521. Navigation pointer into the Stacks Project; the standard fact it supplies is applied inside the counted proof span and the argument resting on it stays there.\par\medskip
\ifdefined\hypertarget\hypertarget{sourcepointer-6}{}\fi
\noindent\textit{Source reference.} Original lemma 0EBN: focused-seeds/AGP/sources/stacks-snapshot-20261011/tags.txt, source lines 17928--17928. Navigation pointer into the Stacks Project; the standard fact it supplies is applied inside the counted proof span and the argument resting on it stays there.\par\medskip
\ifdefined\hypertarget\hypertarget{sourcepointer-7}{}\fi
\noindent\textit{Source reference.} Original lemma 01S8: focused-seeds/AGP/sources/stacks-snapshot-20261011/tags.txt, source lines 2329--2329. Navigation pointer into the Stacks Project; the standard fact it supplies is applied inside the counted proof span and the argument resting on it stays there.\par\medskip
\ifdefined\hypertarget\hypertarget{sourcepointer-8}{}\fi
\noindent\textit{Source reference.} Original lemma 01QF: focused-seeds/AGP/sources/stacks-snapshot-20261011/tags.txt, source lines 2264--2264. Navigation pointer into the Stacks Project; the standard fact it supplies is applied inside the counted proof span and the argument resting on it stays there.\par\medskip
\begin{lemma}
\label{lemma-affine-Xs}
Assumptions and notation as in Lemma \ref{lemma-structure-sheaf-Xs}.
The following are equivalent
\begin{enumerate}
\item the inclusion morphism $j : U \to X$ is affine, and
\item for every $x \in X \setminus U$ there is an $n > 0$
such that $s^n \in \mathfrak m_x \mathcal{F}^{[n]}_x$.
\end{enumerate}
\end{lemma}

\begin{proof}
Assume (1). Then for $x \in X \setminus U$ the inverse image $U_x$ of $U$
under the canonical morphism $f_x : \Spec(\mathcal{O}_{X, x}) \to X$ is affine
and does not contain $x$. Thus $\mathfrak m_x \Gamma(U_x, \mathcal{O}_{U_x})$
is the unit ideal. In particular, we see that we can write
$$
1 = \sum f_i g_i
$$
with $f_i \in \mathfrak m_x$ and $g_i \in \Gamma(U_x, \mathcal{O}_{U_x})$.
By Lemma \ref{lemma-structure-sheaf-Xs} we have
$\Gamma(U_x, \mathcal{O}_{U_x}) = \colim \mathcal{F}^{[n]}_x$
with transition maps given by multiplication by $s$.
Hence for some $n > 0$ we have
$$
s^n = \sum f_i t_i
$$
for some $t_i = s^ng_i \in \mathcal{F}^{[n]}_x$. Thus (2) holds.

\medskip\noindent
Conversely, assume that (2) holds. To prove $j$ is affine is local on $X$,
see Morphisms, Lemma \ref{morphisms-lemma-characterize-affine}.
Thus we may and do assume that $X$ is affine. Our goal is to
show that $U$ is affine.
By Cohomology of Schemes, Lemma \ref{coherent-lemma-affine-if-quasi-affine}
it suffices to show that $H^p(U, \mathcal{O}_U) = 0$ for $p > 0$.
Since $H^p(U, \mathcal{O}_U) = H^0(X, R^pj_*\mathcal{O}_U)$
(Cohomology of Schemes, Lemma
\ref{coherent-lemma-quasi-coherence-higher-direct-images-application})
and since $R^pj_*\mathcal{O}_U$ is quasi-coherent
(Cohomology of Schemes, Lemma
\ref{coherent-lemma-quasi-coherence-higher-direct-images})
it is enough to show the stalk $(R^pj_*\mathcal{O}_U)_x$
at a point $x \in X$ is zero. Consider the base change diagram
$$
\xymatrix{
U_x \ar[d]_{j_x} \ar[r] & U \ar[d]^j \\
\Spec(\mathcal{O}_{X, x}) \ar[r] & X
}
$$
By Cohomology of Schemes, Lemma
\ref{coherent-lemma-flat-base-change-cohomology} we have
$(R^pj_*\mathcal{O}_U)_x = R^pj_{x, *}\mathcal{O}_{U_x}$.
Hence we may assume $X$ is local with closed point $x$
and we have to show $U$ is affine (because this is equivalent to
the desired vanishing by the reference given above).
In particular $d = \dim(X)$ is finite
(Algebra, Proposition \ref{algebra-proposition-dimension}).
If $x \in U$, then $U = X$ and the result is clear.
If $d = 0$ and $x \not \in U$, then $U = \emptyset$
and the result is clear. Now assume $d > 0$ and $x \not \in U$.
Since $j_*\mathcal{O}_U = \colim \mathcal{F}^{[n]}$
our assumption means that we can write
$$
1 = \sum f_i g_i
$$
for some $n > 0$, $f_i \in \mathfrak m_x$, and $g_i \in \mathcal{O}(U)$.
By induction on $d$ we know that $D(f_i) \cap U$ is affine
for all $i$: going through the whole argument just given with
$X$ replaced by $D(f_i)$ we end up with Noetherian local rings
whose dimension is strictly smaller than $d$. Hence $U$
is affine by Properties, Lemma \ref{properties-lemma-characterize-affine}
as desired.
\end{proof}

\end{document}
